% Proposed figure environment. Not inserted into the paper.
% The parent reviews scheme A and scheme B, picks one, then places it
% in the VLDB template. This file is not a standalone paper and is not
% compiled here.
%
% Checked fact: the current paper sources have no figure environment for
% this local pointer rule. Suggested label fig:eg:pointer-rule is new.
% Suggested place: a single-column figure near the unified LT figure.
% Do not reuse fig:eg:loser_tree_kway_group_merge or fig:eg:segment_loser_tree.
%
% PDF media width matches the v2 pointer figure (249.12 pt, about 88 mm).
% Scheme A is 435 pt tall. Scheme B is 617 pt tall. Height grew so the
% tree was not shrunk. \linewidth will scale the chosen PDF.

% --- scheme A: before | after ---
\begin{figure}
  \centering
  \includegraphics[width=\linewidth]{pointer-scheme-a.pdf}
  \caption{Local pointer rule on one internal node ($h=2$).
    Child 0 is run $R_1$ and child 1 is run $R_2$.
    Slot 0 and slot 1 are replay arguments, not child names.
    Build assigns pointer $=$ child 1.
    The next refresh follows that child.
    Refresh returns pointer $=$ child 0; the previous value was child 1.
    A returned pointer may name the same child.
    That case is not in the trace and is not drawn.}
  \label{fig:eg:pointer-rule}
\end{figure}

% --- scheme B: exhaustion, returned pointer, next refresh ---
% Use this environment instead of the one above. Do not include both.
%\begin{figure}
%  \centering
%  \includegraphics[width=\linewidth]{pointer-scheme-b.pdf}
%  \caption{Same local pointer rule as scheme A, drawn as three steps:
%    exhaustion, the returned pointer, then the next refresh.
%    The gap states that the fill after pointer $=$ child 0 is not in the trace.}
%  \label{fig:eg:pointer-rule}
%\end{figure}
